% ---
% title: Two stocks on the same day
% tags: finance, joint-discrete, marginal, joint-distribution-function
% ---

\begin{problem}
	\begin{context}
		The daily price moves of two stocks are usually not unrelated. On a good
		day for the market most stocks rise, and on a bad day most of them fall.
		To describe two moves together, we need their \emph{joint}
		distribution: the probability of every combination of outcomes, not
		only the probabilities for each stock on its own.
	\end{context}
	On a given day, let $X=1$ if stock A rises and $X=-1$ if it falls. Let
	$Y=1$, $0$ or $-1$ if stock B rises, stays unchanged or falls. The joint
	probability mass function $p(x,y)=P(X=x,Y=y)$ is given by the table
	\begin{center}
		\renewcommand{\arraystretch}{1.2}
		\begin{tabular}{c|ccc}
			         & $y=-1$ & $y=0$  & $y=1$  \\
			\hline
			$x=-1$   & $0.20$ & $0.10$ & $0.10$ \\
			$x=1$    & $0.05$ & $0.15$ & $0.40$
		\end{tabular}
	\end{center}
	(a) Find the marginal probability mass functions $p_X$ and $p_Y$.
	(b) Find $P(X=Y)$, the probability that both stocks fall or both rise.
	(c) Find the value of the joint distribution function at $(0,0)$:
	$F(0,0)=P(X\le0,\,Y\le0)$.
	(d) Find the probability that stock B rises, given that stock A rises.
\end{problem}

\begin{slideproblem}
	$X=1$ / $-1$ if stock A rises / falls; $Y=1$, $0$, $-1$ if stock B rises,
	stays, falls. The joint pmf $p(x,y)=P(X=x,Y=y)$ is in the table below.
	(a) Marginal pmfs $p_X$, $p_Y$. \ (b)~$P(X=Y)$. \
	(c)~$F(0,0)=P(X\le0,\,Y\le0)$. \ (d)~$P(\text{B rises}\mid\text{A rises})$.
\end{slideproblem}

\begin{slides}
	\small
	\exstatement
	\begin{center}
		\renewcommand{\arraystretch}{1.25}
		\begin{tabular}{c|ccc|c}
			       & $y=-1$ & $y=0$  & $y=1$  & $p_X(x)$ \\
			\hline
			$x=-1$ & $0.20$ & $0.10$ & $0.10$ & \uncover<3->{$\alert{0.4}$} \\
			$x=1$  & $0.05$ & $0.15$ & $0.40$ & \uncover<3->{$\alert{0.6}$} \\
			\hline
			$p_Y(y)$ & \uncover<4->{$\alert{0.25}$} & \uncover<4->{$\alert{0.25}$}
			& \uncover<4->{$\alert{0.5}$} & \uncover<4->{$1$}
		\end{tabular}
	\end{center}
	\uncover<2->{\sol (a) Marginals: add up the rows and the columns; write
		the sums in the \alert{margins}.}\\
	\uncover<3->{Row sums: $p_X(-1)=0.2+0.1+0.1=0.4$, \ $p_X(1)=0.05+0.15+0.4=0.6$.}\\
	\uncover<4->{Column sums: $p_Y(-1)=0.25$, \ $p_Y(0)=0.25$, \ $p_Y(1)=0.5$.}
\end{slides}

\begin{slides}
	\small
	(b) $X=Y$: both fall $(-1,-1)$ or both rise $(1,1)$. Disjoint cells: add.
	\[ P(X=Y)=p(-1,-1)+p(1,1)=0.20+0.40=\alert{0.60}. \]
	\pause
	(c) $X\le0$ means $X=-1$; \ $Y\le0$ means $Y=-1$ or $Y=0$:
	\[ F(0,0)=P(X\le0,\,Y\le0)=p(-1,-1)+p(-1,0)=0.20+0.10=\alert{0.30}. \]
	\pause
	(d) Conditional probability: the joint cell divided by the marginal.
	\[ P(Y=1\mid X=1)=\frac{p(1,1)}{p_X(1)}=\frac{0.40}{0.60}=\frac23
		\approx\alert{0.667}. \]
	\pause
	Compare: $P(Y=1)=p_Y(1)=0.5$. Knowing that A rises makes it more likely
	that B rises too.
\end{slides}

\begin{solution}
	\idea All the information is in the joint table. A marginal pmf is a row
	or column sum of the table. Any other probability is the sum of the cells
	that belong to the event, because different cells are mutually exclusive
	outcomes.

	\step{(a)} The event $\{X=-1\}$ happens together with exactly one of the
	values $Y=-1$, $Y=0$, $Y=1$. These three events are mutually exclusive,
	so
	\[ p_X(-1)=P(X=-1)=p(-1,-1)+p(-1,0)+p(-1,1)=0.20+0.10+0.10=0.4. \]
	In the same way, the second row gives
	\[ p_X(1)=p(1,-1)+p(1,0)+p(1,1)=0.05+0.15+0.40=0.6. \]
	For $Y$ we add the columns instead:
	\[ p_Y(-1)=0.20+0.05=0.25,\qquad p_Y(0)=0.10+0.15=0.25,\qquad
		p_Y(1)=0.10+0.40=0.5. \]
	It is natural to write these sums in the margins of the table. This is
	where the name \emph{marginal} distribution comes from:
	\begin{center}
		\renewcommand{\arraystretch}{1.25}
		\begin{tabular}{c|ccc|c}
			         & $y=-1$ & $y=0$  & $y=1$  & $p_X(x)$ \\
			\hline
			$x=-1$   & $0.20$ & $0.10$ & $0.10$ & $0.4$ \\
			$x=1$    & $0.05$ & $0.15$ & $0.40$ & $0.6$ \\
			\hline
			$p_Y(y)$ & $0.25$ & $0.25$ & $0.5$  & $1$
		\end{tabular}
	\end{center}
	Check: both margins add up to $1$ ($0.4+0.6=1$ and $0.25+0.25+0.5=1$),
	as every pmf must. So stock A rises with probability $0.6$, and stock B
	rises with probability $0.5$.

	\step{(b)} $X$ takes only the values $-1$ and $1$, so $X=Y$ can happen in
	two ways: both stocks fall, $(X,Y)=(-1,-1)$, or both rise,
	$(X,Y)=(1,1)$. The value $Y=0$ can never equal $X$. The two cells are
	mutually exclusive, so we add them:
	\[ P(X=Y)=p(-1,-1)+p(1,1)=0.20+0.40=0.60. \]
	On 60\% of the days the two stocks move in the same direction.

	\step{(c)} The joint distribution function is
	$F(a,b)=P(X\le a,\,Y\le b)$: both conditions must hold. At $(0,0)$:
	$X\le0$ holds only for $X=-1$, and $Y\le0$ holds for $Y=-1$ and $Y=0$. So
	we add the cells in the first row and the first two columns:
	\[ F(0,0)=P(X\le0,\,Y\le0)=p(-1,-1)+p(-1,0)=0.20+0.10=0.30. \]
	In words: with probability $0.30$, stock A falls and stock B does not
	rise.

	\step{(d)} We need $P(Y=1\mid X=1)$. By the definition of conditional
	probability, it is the probability of both events divided by the
	probability of the condition. The numerator is one cell of the table,
	the denominator is the marginal from (a):
	\[ P(Y=1\mid X=1)=\frac{P(X=1,\,Y=1)}{P(X=1)}=\frac{p(1,1)}{p_X(1)}
		=\frac{0.40}{0.60}=\frac23\approx0.667. \]
	Without any information about stock A, stock B rises with probability
	$p_Y(1)=0.5$. Knowing that A rises raises this to $2/3$: on a day when A
	rises, B tends to rise as well.
	\begin{answer}
		(a) $p_X(-1)=0.4$, $p_X(1)=0.6$; \ $p_Y(-1)=0.25$, $p_Y(0)=0.25$,
		$p_Y(1)=0.5$
		\quad (b) $P(X=Y)=0.60$
		\quad (c) $F(0,0)=0.30$
		\quad (d) $P(Y=1\mid X=1)=\frac23\approx0.667$
	\end{answer}
\end{solution}
